\scanpage{116}
% archive sheet number: 115
et on introduit sur l'ens \inserted{$\coprod_{i\in\{0,1,2\}}\Phi_i$} des figures universelles une relation d'ordre ou d'incidence par
\[
 x\leq y\ \Longleftrightarrow\
 \left\{\begin{aligned}
 \text{a)}\ &\dim(x)\leq\dim(y)\\
 \text{b)}\ &\exists z\in\Phi_*,\ \text{avec }x,y\unsure{\preceq|\leq}z
 \end{aligned}\right.\tag*{(9)}
\]

Si $\deleted{\unsure{E}}X$ est un \deleted{\unsure{module}} \deleted{\unsure{fibré vectoriel}} affine ou \deleted{projectif} \inserted{de dim \deleted{\unsure{rel}} $>3$} \deleted{\unsure{sur une}} sur une base quelconque, on définit les drapeaux de type $H$ \inserted{$(H\subset\{0,1,2\},\ H\ne\varnothing)$} de $X$ de façon habituelle, \deleted{Une \unsure{représentation|réalisation}} $\Phi_H(X/S)$, $\Phi_{\unsure{**|\bullet *}}(X/S)$, $\Phi_*(X/S)$ \ldots
\margin{\unsure{Quasi simplt|ou simplt}}

Un polyèdre \uline{régulier} \inserted{$P$} \inserted{\deleted{\unsure{splitté}} scindé} dans $X$ est la donnée, \inserted{d'abord,} de trois ensembles \inserted{(finis)}
% slanted
\margin{La condition de finitude n'est \unsure{peut-être} pas essentielle\unsure{\ldots}}
\[
 \left\{\begin{aligned}
 S:=S_P&\subset\Phi_0(X/S)=\Gamma(X/S)\qquad\text{ens des sommets}\\
 A:=A_P&\subset\Phi_1(X/S)\deleted{=\unsure{A}}\qquad\text{\unsure{..} arêtes}\\
 F:=F_P&\subset\Phi_2(X/S)\qquad\text{\unsure{..} faces}
 \end{aligned}\right.\tag*{(10)}
\]
satisfaisant les conditions
\[
 \left\{\begin{aligned}
 \text{(a)}\quad &\forall a\in A,\ \text{\deleted{il y a \unsure{deux exactement}} \inserted{l'ens. des}}\ s\in S\ \text{incidents}\\
 &\text{à }\unsure{S|a}\text{ est de cardinal 2 \unsure{exactt|exactement} --- et ceci reste}\\
 &\text{vrai après }\text{tt chgt de base }S'\to S,\ S'\text{ non vide}\\
 \text{(b)}\quad &\forall s\in S,\ f\in F,\ \text{l'ens des }a\text{ incidentes}\\
 &\text{à }s\text{ et }f\text{ est \inserted{loc} de cardinal égal à }\substack{2\\0}\text{ --- et}\\
 &\text{ceci reste vrai après }\text{tt chgt de base }S'\to S,\ S'\ne\varnothing\\
 \text{(c)}\quad &\forall a\in A,\ \text{l'ens des }f\in F\text{ incidentes à }a\\
 &\text{est de cardinal }=\unsure{2|?},\ \text{et ceci reste vrai}\\
 &\text{après }\text{tt chgt. de base }S'\to S,\ S'\ne\varnothing
 \end{aligned}\right.\tag*{\boxed{11}}
\]
